\appendix
\section*{Annexes}

\begin{frame}{Annexe --- Architecture détaillée}
  \scriptsize
  \centering
  \begin{tikzpicture}[
      card/.style={draw=McmTerracotta,rounded corners=2pt,fill=white,
                   minimum height=1.10cm,text width=1.85cm,align=center,
                   inner sep=3pt,font=\scriptsize},
      ext/.style={card,draw=McmMuted,dashed},
      flow/.style={{Latex[length=2mm]}-{Latex[length=2mm]},thick,
                   draw=McmTerracotta},
      lab/.style={font=\tiny,text=McmMuted,inner sep=1pt}]
    \node[ext]  (user)  at (0,0)    {\textbf{Navigateur}\\PC ou mobile};
    \node[card] (front) at (3.2,0)  {\textbf{Front}\\SvelteKit};
    \node[card] (api)   at (6.4,0)  {\textbf{API}\\FastAPI};
    \node[card] (db)    at (9.6,0)  {\textbf{Base}\\PostgreSQL};
    \node[ext]  (ai)    at (6.4,-1.75) {\textbf{Modules IA}\\prédiction,\\ordonnanceur};
    \draw[flow] (user)  -- node[lab,above=2pt] {HTTPS} (front);
    \draw[flow] (front) -- node[lab,above=2pt] {interne} (api);
    \draw[flow] (api)   -- node[lab,above=2pt] {SQL} (db);
    \draw[flow] (api)   -- node[lab,right=2pt] {contrats Python} (ai);
    \begin{scope}[on background layer]
      \node[draw=McmSage,dashed,rounded corners=3pt,inner sep=4pt,
            fit=(front)(api)(db)(ai)] (srv) {};
    \end{scope}
    \node[font=\tiny,text=McmSage,anchor=south east]
      at ([shift={(0,1pt)}]srv.north east) {Podman rootless, Ansible};
  \end{tikzpicture}

  \vspace{0.12cm}
  \begin{columns}[T,onlytextwidth]
    \begin{column}{0.49\textwidth}
      \footnotesize Le navigateur ne parle qu'au front ; l'API garde les
      droits et les données.
    \end{column}
    \begin{column}{0.49\textwidth}
      \footnotesize Argon2, cookies HttpOnly/Secure, TLS, conteneurs
      rootless, données pseudonymisées.
    \end{column}
  \end{columns}
\end{frame}

\begin{frame}{Annexe --- Variables, contraintes et objectif}
  \scriptsize
  \begin{columns}[T,onlytextwidth]
    \begin{column}{0.49\textwidth}
      \begin{block}{Décision}
        \begin{tabular}{@{}ll@{}}
          $P=\{1,\dots,n\}$ & patients \\
          $V=\{1,\dots,m\}$ & vacations \\
          $x_{i,v}\in\{0,1\}$ & patient $i$ dans la vacation $v$ \\
        \end{tabular}
        \[
          \sum_{v\in V} x_{i,v}=1 \quad \forall i\in P
        \]
      \end{block}
    \end{column}

    \begin{column}{0.49\textwidth}
      \begin{block}{Contraintes dures}
        \begin{itemize}
          \item capacité : $\sum_i d_i x_{i,v}\le C_v$ ;
          \item lits : $\sum_i S_i x_{i,v}\le L_j$ pour chaque jour $j$ ;
          \item fenêtre d'urgence : $x_{i,v}=0$ hors de la fenêtre ;
          \item spécialité : vacation compatible avec le patient.
        \end{itemize}
      \end{block}
    \end{column}
  \end{columns}

  \vspace{0.12cm}
  \begin{alertblock}{Objectif réellement implémenté}
    \[
      \text{coût}=5\cdot\text{dépassement}_{vac}
      +3\cdot\text{dépassement}_{lits}
      +0{,}05\cdot\text{écart-type}_{vac}
    \]
    La \textbf{fitness} vaut $-\text{coût}$ : on maximise, $0$ correspond à un
    planning faisable et équilibré.
  \end{alertblock}
  \source{Durées $d_i$ et séjours $S_i$ estimés par les modèles ; capacité
    de vacation typique : 240 min.}
\end{frame}

\begin{frame}{Annexe --- Les dix méthodes}
  \scriptsize
  \renewcommand{\arraystretch}{1.10}
  \begin{tabular}{@{}>{\raggedright\arraybackslash}p{3.0cm}>{\raggedright\arraybackslash}p{5.4cm}>{\raggedright\arraybackslash}p{4.2cm}@{}}
    \toprule
    \textbf{Méthode} & \textbf{Principe} & \textbf{Rôle} \\
    \midrule
    Recuit simulé & accepte une dégradation avec une probabilité décroissante
      & exploration large \\
    Tabou & meilleur voisin non tabou, mémoire des mouvements
      & intensification locale \\
    Génétique & population, sélection, croisement, mutation
      & diversité des solutions \\
    Tabou $\times$ Recuit & voisinage tabou et acceptation type recuit
      & hybride local \\
    Fourmis (ACO) & construction par phéromones
      & exploration par construction \\
    Génétique $\times$ Tabou & tabou sur le meilleur individu de chaque
      génération & hybride \\
    Génétique $\times$ Recuit & recuit court sur le meilleur individu
      & hybride \\
    Fourmis $\times$ Tabou & chaque fourmi est intensifiée par tabou
      & hybridation \\
    SMA & agents simples, tableau noir, coordinateur
      & coordination \\
    SMA $\times$ Méta & agents recuit, tabou, génétique, fourmis
      & coordination métaheuristique \\
    \bottomrule
  \end{tabular}
  \source{Chaque méthode partage la même interface et le même budget temps.}
\end{frame}

\begin{frame}{Annexe --- Benchmark 300 patients, 80 vacations}
  \scriptsize
  \begin{columns}[T,onlytextwidth]
    \begin{column}{0.56\textwidth}
      \renewcommand{\arraystretch}{1.02}
      \begin{tabular}{@{}lrrr@{}}
        \toprule
        Méthode & Écart & Rang & IQR \\
        \midrule
        Recuit simulé & \textbf{0,015} & \textbf{1,25} & 0,008 \\
        Tabou $\times$ Recuit & \textbf{0,018} & 1,75 & 0,013 \\
        Tabou & 0,240 & 3,55 & 0,070 \\
        Génétique $\times$ Tabou & 0,261 & 3,75 & 0,106 \\
        SMA & 1,078 & 5,05 & 0,194 \\
        SMA $\times$ Méta & 1,196 & 6,15 & 0,154 \\
        Génétique $\times$ Recuit & 2,031 & 6,50 & 0,323 \\
        Fourmis $\times$ Tabou & 1\,006,0 & 8,00 & 180,9 \\
        Fourmis (ACO) & 2\,031,2 & 9,00 & 256,1 \\
        Génétique & 2\,836,2 & 10,00 & 380,1 \\
        \bottomrule
      \end{tabular}
    \end{column}
    \begin{column}{0.42\textwidth}
      \centering
      \includegraphics[width=\linewidth,height=0.34\textheight,keepaspectratio]{grande_convergence.png}\\[2pt]
      {\tiny Convergence médiane, 20 graines, 3 s.}\\[4pt]
      \footnotesize Recuit et Tabou $\times$ Recuit restent sous un écart de
      0,02 pour une référence de $-38\,847\,195$.
    \end{column}
  \end{columns}
  \source{20 graines, budget 3 s, référence à la meilleure solution observée.}
\end{frame}

\begin{frame}{Annexe --- Échelle réelle : 582 patients, 280 vacations}
  \scriptsize
  \begin{columns}[T,onlytextwidth]
    \begin{column}{0.50\textwidth}
      \renewcommand{\arraystretch}{1.12}
      \begin{tabular}{@{}lrr@{}}
        \toprule
        Méthode & Médiane & Rang moyen \\
        \midrule
        Recuit simulé & $-16{,}3$ & 1,85 \\
        Tabou $\times$ Recuit & $-14{,}9$ & 2,05 \\
        Tabou & $-16{,}4$ & 2,65 \\
        Génétique $\times$ Tabou & $-23{,}9$ & 3,50 \\
        SMA & $-91{,}8$ & 5,85 \\
        Fourmis (ACO) & $-3\,438{,}8$ & 9,70 \\
        \bottomrule
      \end{tabular}\\[4pt]
      {\tiny 20 graines ; référence $-8\,783$.}
    \end{column}
    \begin{column}{0.48\textwidth}
      \centering
      \includegraphics[width=\linewidth,height=0.33\textheight,keepaspectratio]{reelle_convergence.png}\\[2pt]
      \footnotesize Le couloir faisable est étroit : la meilleure solution
      occupe 42 lits sur 42.
    \end{column}
  \end{columns}
  \source{Parquet EDA prétraité, 40 jours, capacité 600 min par vacation,
    4 s par méthode et par graine.}
\end{frame}

\begin{frame}{Annexe --- Quatre aléas et quatre profils de repli}
  \scriptsize
  \renewcommand{\arraystretch}{1.08}
  \renewcommand{\arraystretch}{1.18}
  \begin{tabular}{@{}>{\raggedright\arraybackslash}p{2.6cm}>{\raggedright\arraybackslash}p{5.4cm}>{\raggedright\arraybackslash}p{4.6cm}@{}}
    \toprule
    \textbf{Aléa} & \textbf{Description} & \textbf{Impact} \\
    \midrule
    Urgence & patient imprévu, fenêtre imposée
      & insertion dans une vacation compatible \\
    Annulation & patient programmé qui annule
      & libère un créneau et des lits \\
    Indisponibilité lits & capacité réduite sur un intervalle
      & réaffecter les séjours chevauchants \\
    Retard bloc & prolongation ou incident
      & capacité utile réduite \\
    \bottomrule
  \end{tabular}

  \vspace{0.15cm}
  \begin{columns}[T,onlytextwidth]
    \begin{column}{0.49\textwidth}
      \begin{block}{Profils proactifs}
        \begin{itemize}
          \item nominal : 100\,\% de la capacité ;
          \item robuste : 15\,\% de marge bloc et lits ;
          \item sécurité lits : lissage strict ;
          \item Date B : bonne qualité, différente du nominal.
        \end{itemize}
      \end{block}
    \end{column}
    \begin{column}{0.49\textwidth}
      \begin{exampleblock}{Double proposition}
        \footnotesize Pour chaque patient, l'outil produit une Date A
        optimale et une Date B alternative, avec la raison de l'écart.
      \end{exampleblock}
    \end{column}
  \end{columns}
\end{frame}

\begin{frame}{Annexe --- Adaptation réactive et arbitrage}
  \scriptsize
  \begin{columns}[T,onlytextwidth]
    \begin{column}{0.52\textwidth}
      \begin{block}{Score}
        \[
          \begin{aligned}
            \text{Score} = \text{Fitness}
            &- w_p \cdot \text{déplacements}\\
            &- w_j \cdot \text{changements de jour}
          \end{aligned}
        \]
        Le passé est gelé ; les urgences entrent dans leur fenêtre ; les
        patients futurs bougent le moins possible.
      \end{block}
    \end{column}
    \begin{column}{0.46\textwidth}
      \begin{exampleblock}{Rapport d'arbitrage}
        \footnotesize Pour chaque déplacement : patient, ancienne date,
        nouvelle date, motif. Pour chaque urgence : vacation retenue et
        patients décalés.
      \end{exampleblock}
    \end{column}
  \end{columns}

  \vspace{0.15cm}
  \centering
  \renewcommand{\arraystretch}{1.18}
  \begin{tabular}{llrr}
    \toprule
    Situation & Politique & Terminées / 279 & Dépassement (min) \\
    \midrule
    Sans fermeture & Figée ou réactive & 262 & 2\,954 \\
    Avec fermeture & Figée & \textbf{255} & 2\,911 \\
    Avec fermeture & Réactive & 251 & \textbf{2\,716} \\
    \bottomrule
  \end{tabular}
  \source{42 lits, 2 salles, 3--31 janvier 2022, graine 0. La replanification
    réduit le dépassement de 195 min mais termine 4 interventions de moins.}
\end{frame}

\begin{frame}{Annexe --- Simulations dynamiques}
  \scriptsize
  \centering
  \begin{tikzpicture}[x=1.0cm,y=1cm,>=Latex]
    \draw[very thick,draw=McmInk] (0,0) -- (10.3,0);
    \draw[line width=4pt,draw=McmAmber] (2.3,0) -- (5.5,0);
    \foreach \x/\lab in {0/08h,2.3/10h,5.5/12h,10.3/17h} {
      \draw[draw=McmInk] (\x,-0.12) -- (\x,0.12);
      \node[below=5pt,font=\scriptsize] at (\x,0) {\lab};
    }
    \node[above=8pt,text=McmAmber,font=\small\bfseries] at (3.9,0)
      {Salle 1 : aucune nouvelle entrée};
  \end{tikzpicture}

  \vspace{0.25cm}
  \begin{columns}[T,onlytextwidth]
    \begin{column}{0.49\textwidth}
      \begin{block}{Plan figé}
        \footnotesize Les affectations et l'ordre initial sont conservés.
        Les cas suivants attendent la réouverture.
      \end{block}
    \end{column}
    \begin{column}{0.49\textwidth}
      \begin{exampleblock}{Plan réactif}
        \footnotesize À 10h et à 12h, le coordinateur soumet un nouveau
        planning pour les cas en attente.
      \end{exampleblock}
    \end{column}
  \end{columns}

  \vspace{0.10cm}
  \centering\footnotesize Un acte commencé se termine. Une proposition
  invalide est rejetée : l'ancien plan reste actif.
  \source{Mesa 3.5 ; urgences, annulations et fermetures de lits restent à
    exécuter dans ce simulateur.}
\end{frame}

\begin{frame}{Annexe --- Données EDA et nettoyage}
  \scriptsize
  \begin{columns}[T,onlytextwidth]
    \begin{column}{0.31\textwidth}
      \begin{exampleblock}{Volume}
        \centering {\Large\bfseries 14\,649}\\ cas hospitaliers
      \end{exampleblock}
    \end{column}
    \begin{column}{0.31\textwidth}
      \begin{block}{Patients}
        \centering {\Large\bfseries 11\,108}\\ identifiants distincts
      \end{block}
    \end{column}
    \begin{column}{0.31\textwidth}
      \begin{alertblock}{Variables}
        \centering {\Large\bfseries 30}\\ administratives, cliniques, bloc
      \end{alertblock}
    \end{column}
  \end{columns}

  \vspace{0.15cm}
  \begin{columns}[T,onlytextwidth]
    \begin{column}{0.49\textwidth}
      \begin{block}{Anomalies}
        \begin{itemize}
          \item 2\,560 incohérences de durée de séjour, soit 17,5\,\% ;
          \item 1\,884 séjours notés « 1 jour » malgré une nuit ;
          \item 552 incisions à zéro ; 42 incisions avant l'entrée.
        \end{itemize}
      \end{block}
    \end{column}
    \begin{column}{0.49\textwidth}
      \begin{exampleblock}{Nettoyage}
        \begin{itemize}
          \item 142 lignes sans type d'intervention ou praticien ;
          \item 14\,507 lignes conservées ;
          \item 2\,488 durées corrigées, durée source gardée.
        \end{itemize}
      \end{exampleblock}
    \end{column}
  \end{columns}
  \source{La granularité est le séjour, pas le patient ; 2\,546 patients
    apparaissent dans plusieurs cas.}
\end{frame}

\begin{frame}{Annexe --- KYST : sécurité et confidentialité}
  \scriptsize
  \begin{columns}[T,onlytextwidth]
    \begin{column}{0.32\textwidth}
      \begin{block}{Comptes et sessions}
        \begin{itemize}
          \item rôles vérifiés par l'API ;
          \item mots de passe hachés avec Argon2 ;
          \item cookies HttpOnly et Secure ;
          \item formulaires protégés contre le CSRF.
        \end{itemize}
      \end{block}
    \end{column}
    \begin{column}{0.32\textwidth}
      \begin{exampleblock}{Données patient}
        \begin{itemize}
          \item pseudonymisées ;
          \item tableaux de bord agrégés ;
          \item auteur et date de chaque décision ;
          \item aucune donnée réelle dans le dépôt.
        \end{itemize}
      \end{exampleblock}
    \end{column}
    \begin{column}{0.32\textwidth}
      \begin{alertblock}{Infrastructure}
        \begin{itemize}
          \item API et base non exposées ;
          \item seul le serveur web sécurisé joint le front ;
          \item conteneurs sans privilèges ;
          \item secrets hors du dépôt.
        \end{itemize}
      \end{alertblock}
    \end{column}
  \end{columns}
  \source{Prévu : journal d'audit complet et droits limités aux demandes de
    chaque chirurgien.}
\end{frame}

\begin{frame}{Annexe --- Chiffres et sources}
  \scriptsize
  \renewcommand{\arraystretch}{1.12}
  \begin{tabular}{@{}>{\raggedright\arraybackslash}p{1.4cm}>{\raggedright\arraybackslash}p{7.9cm}>{\raggedright\arraybackslash}p{2.7cm}@{}}
    \toprule
    \textbf{Chiffre} & \textbf{Lecture} & \textbf{Source} \\
    \midrule
    64,4\,\% & chirurgie ambulatoire en France en 2023 & HAS, IQSS 2026 \\
    70\,\% & objectif national fixé pour 2022 & HAS, IQSS 2026 \\
    80\,\% & cap affiché pour la suite & HAS, IQSS 2026 \\
    85\,\% & cible d'occupation des salles & ANAP \\
    $\sim$10\,\% & interventions programmées annulées, souvent reprogrammées
      & HAS, Flash 2026 \\
    90\,\% & établissements avec un bloc commun & HAS, IQSS 2026 \\
    86\,\% & personnel partagé avec l'hospitalisation & HAS, IQSS 2026 \\
    35\,\% & anticipation de la douleur & HAS, IQSS 2026 \\
    75\,\% & contact patient entre J1 et J3 & HAS, IQSS 2026 \\
    \bottomrule
  \end{tabular}
  \source{HAS, \emph{Organisation en unité de chirurgie ambulatoire et IQSS},
    avril 2026 ; HAS, \emph{Reprogrammation au bloc opératoire}, Flash
    Sécurité Patient, 2026 ; ANAP, \emph{Outil d'analyse bloc opératoire}.}
\end{frame}
